\documentclass[11pt]{article}
\usepackage[T1]{fontenc}
\usepackage{lmodern,microtype,amsmath,amssymb,amsthm,mathtools}
\usepackage[margin=1in]{geometry}
\usepackage{booktabs,array,longtable,graphicx,xcolor,enumitem}
\usepackage[colorlinks=true,linkcolor=blue!45!black,urlcolor=blue!55!black,citecolor=blue!45!black]{hyperref}
\hypersetup{pdftitle={Causal memory for the homogeneous tracial partial-SWAP stream: proved bounds and unresolved tradeoff},pdfauthor={Constructed research return},pdfsubject={General causal receiver converse, physical compression obstructions, and exact finite checks}}
\setlength{\parindent}{0pt}
\setlength{\parskip}{4pt}
\setlist{nosep,leftmargin=*}
\newtheorem{theorem}{Theorem}
\newtheorem{lemma}[theorem]{Lemma}
\newtheorem{corollary}[theorem]{Corollary}
\newcommand{\Tr}{\operatorname{Tr}}
\newcommand{\id}{\operatorname{id}}
\newcommand{\Ad}{\operatorname{Ad}}
\newcommand{\rank}{\operatorname{rank}}
\newcommand{\Span}{\operatorname{span}}
\newcommand{\norm}[1]{\left\lVert#1\right\rVert}
\newcommand{\dd}{d_{\diamond}}
\title{Causal memory for the homogeneous\newline tracial partial-SWAP stream}
\author{Proved bounds, physically specified candidates, and the unresolved tradeoff}
\date{5 November 2026}
\begin{document}
\maketitle
\begin{abstract}
The source is fixed: one arbitrary qubit, possibly entangled with a reference,
interacts once with each independent maximally mixed carrier through the same
known weak partial SWAP. The source is discarded. No bath purifiers, preparation
keys, or return link are available. The accepted all-output decoder acquires the
payload at the quartic collision scale, but retains the entire stream.

This return does \emph{not} establish the requested optimal causal
memory/processing tradeoff. Its principal completed result is a converse for
every permitted causal receiver: successful payload recovery forces a
bistochastic low-Schmidt-number Choi state to have small energy for an exact
bath-scattering singular operator. This yields finite-error spectral-profile,
mode-count, heat-trace, branch-tail, and expected-occupancy bounds. The needed
large-prefix spectral estimate is not proved.

Several concrete compression proposals are resolved at their actual scopes.
Ordinary measured-label Schur compression has an asymptotic half-diamond error
floor of one half. Classical consecutive-pair fusion loses information at the
same fourth order at which it is acquired. A two-qubit folded reverse-decoder
statistic violates causal distinguishability monotonicity. Even reverse-decoded
blocks retaining both a dummy qubit and its locally prepared Bell purifier need
blocks on the original quartic scale, despite arbitrary collective final
decoding. Exact proofs, finite matrix checks, and adaptive-control diagnostics
are supplied. None of these restricted obstructions is substituted for a
general impossibility theorem.
\end{abstract}

\section{The fixed operation and the status of the question}

Let $S$ be the original source qubit and $E_1,\ldots,E_N$ independent carriers
in $\tau=I/2$. Write
\[
 U_j=cI+is\,\mathrm{SWAP}_{S,E_j},\qquad
 c=\cos\theta,\quad s=\sin\theta,\qquad
 U^{(N)}=U_N\cdots U_1.
\]
The acquired channel is
\[
 \mathcal C_N(\rho)=\Tr_S U^{(N)}(\rho\otimes\tau^{\otimes N})U^{(N)\dagger}.
\]
A causal receiver can act on each arriving carrier, keep quantum workspace,
measure, retain unlimited classical records, use feedback, and process retained
blocks. Every coherent delay line, unprocessed carrier, protected register, and
initialized ancilla crossing a time cut counts as quantum memory. It cannot
interact with $S$ again. The requested guarantee is
\[
 \dd(\mathcal D\mathcal C_N,\id)
 :=\tfrac12\norm{\mathcal D\mathcal C_N-\id}_{\diamond}\le\epsilon.
\]
The reference in the diamond norm is essential; estimation of an input Bloch
vector is not the operation.

The accepted all-output construction prepares an independent tracial dummy,
applies the inverse collisions in reverse chronological order, and returns the
dummy. With $r=1-s^4$, its exact channel is
\begin{equation}
 \mathcal D_N\mathcal C_N=(1-r^N)\id+r^N\mathcal R_{I/2},
 \qquad \dd(\mathcal D_N\mathcal C_N,\id)=\tfrac34r^N.
 \label{eq:baseline}
\end{equation}
The accepted arbitrary-decoder converse is
\begin{equation}
 \epsilon_{\rm opt}\ge\tfrac34(1-\sqrt{1-r^N}).
 \label{eq:baseline-lower}
\end{equation}
Thus an attainable collision count is
\begin{equation}
 N_{\rm all}(\theta,\epsilon)=
 \left\lceil\frac{\log(3/(4\epsilon))}{-\log(1-\sin^4\theta)}\right\rceil
 =\Theta\!\left(\theta^{-4}\log(1/\epsilon)\right)
 \label{eq:N}
\end{equation}
in the weak-angle regime. This construction uses $N$ retained output qubits,
one dummy, and $N$ inverse two-qubit gates before routing and compilation.
At $\theta=0.1$, $\epsilon=0.01$, it has $N=43\,462$ and ideal error
$0.0099996156811$. The accepted converse requires at least $36\,552$ collisions.

The task left unresolved here is whether a general causal receiver can attain
\eqref{eq:N} with quantum storage $o(N)$ and near-linear total processing,
jointly in angle, accuracy, and stream length. A smaller register at fixed
angle obtained by hiding an inverse power of $\theta$ or $\epsilon$ is not an
answer. Nor is a successful reader for a programmed late-selector emitter.

\section{A converse for general causal receivers}

This result uses the actual homogeneous source. It does not assume a Schur
update, a tracial receiver, a unitary receiver, a fixed memory dimension, or
that the receiver reconstructs all emitted outputs.

\subsection{The physical cut and an exact operator}

At a prefix of length $\ell$, put $D=2^\ell$, $B=E_1\cdots E_\ell$, and
\[
 J_a=U^{(\ell)}(\sigma_a\otimes I_B)U^{(\ell)\dagger},\qquad a=x,y,z.
\]
All processing before the cut is an instrument
\[
 \mathcal A:B\longrightarrow\bigoplus_y M_y,
\]
where $y$ is classical and $m_y=\dim M_y$ is the actual branch capacity.
For a lower bound, grant the decoder the actual source $S$ at the cut. It can
then simulate all subsequent known source collisions with fresh tracial
carriers and run the original continuation. Thus a successful original
receiver also recovers from
$(\id_S\otimes\mathcal A)\Ad_{U^{(\ell)}}(\rho\otimes I_B/D)$ with the
same error. Granting $S$ is a proof relaxation, not an access resource in the
claimed physical task.

Define the bath-scattering channel
\[
 T_\ell(K)=\tfrac12\Tr_S U^{(\ell)}(I_S\otimes K)U^{(\ell)\dagger}
\]
and its Hilbert--Schmidt adjoint $T_\ell^*$. Both are unital and trace preserving.
On the $D^2$-dimensional bath-operator space set
\begin{equation}
 H_\ell=I-T_\ell T_\ell^*,\qquad
 0=\lambda_0\le\lambda_1\le\cdots\le\lambda_{D^2-1}\le1.
 \label{eq:H}
\end{equation}
No unproved uniqueness of the zero mode is needed. The identity is one zero
mode; exceptional angles can produce others.

The elementary Pauli twirl gives the exact identities
\begin{align}
 \frac16\sum_a\Tr_S[J_a(I\otimes K)J_a]
   &=\frac{4T_\ell T_\ell^*(K)-K}{3},\\
 \frac13\sum_a\norm{[I\otimes K,J_a]}_2^2
   &=\frac{16}{3}\langle K,H_\ell K\rangle.
 \label{eq:dirichlet}
\end{align}
Here $\norm{\cdot}_2$ and the inner product are unnormalized Hilbert--Schmidt
quantities. The ordering $T_\ell T_\ell^*$ is part of the definition.

\subsection{Common recovery witness}

Let $\omega=\mathcal A(I_B/D)$ and define the bath Petz map on its support by
\[
 \mathcal P(X)=\frac1D\mathcal A^*(\omega^{-1/2}X\omega^{-1/2}),
 \qquad W=\mathcal P\mathcal A.
\]
Extend the map trace preservingly on unused support. $W$ is unital and trace
preserving and factors through the same hybrid memory. It is an analytic
witness; the physical receiver is not asked to apply it.

For the two logical eigenstates in Pauli direction $a$, successful recovery
implies
\[
 t_a=\norm{(\id_S\otimes\mathcal A)(J_a)/(2D)}_1\ge1-2\epsilon.
\]
Write $\widetilde\omega=I_S/2\otimes\omega$ and
$x_a=(\id_S\otimes\mathcal A)(J_a)/(2D)$. Weighted Cauchy--Schwarz yields
\begin{equation}
 q_a:=\frac{\Tr[J_a(\id_S\otimes W)(J_a)]}{2D}
 =\Tr[x_a\widetilde\omega^{-1/2}x_a\widetilde\omega^{-1/2}]
 \ge t_a^2.
 \label{eq:q}
\end{equation}
To see the last inequality, pair
$\widetilde\omega^{-1/4}x_a\widetilde\omega^{-1/4}$ with
$\widetilde\omega^{1/4}\operatorname{sign}(x_a)\widetilde\omega^{1/4}$.
The squared norm of the second factor is at most one: in an eigenbasis of
$\widetilde\omega$, use $\sqrt{w_iw_j}\le(w_i+w_j)/2$ and
$\operatorname{sign}(x_a)^2\le I$.

If $K_\alpha$ are Kraus matrices for $W$, its normalized Choi state is
\[
 Z=\frac1D\sum_\alpha|K_\alpha\rangle\!\rangle
                      \langle\!\langle K_\alpha|.
\]
Both marginals equal $I/D$. Equations \eqref{eq:dirichlet}--\eqref{eq:q}
give exactly
\begin{equation}
 E:=\Tr(H_\ell Z)=\frac34\left(1-\frac{q_x+q_y+q_z}{3}\right)
 \le3\epsilon(1-\epsilon)=:e_\epsilon.
 \label{eq:energy}
\end{equation}
The conversion from payload fidelity to this small-energy Choi state is the
central general statement.

\subsection{Capacity, spectral profile, and occupancy}

\begin{theorem}[General prefix memory bound]
Suppose a causal receiver attains half-diamond error $0\le\epsilon\le1/2$,
and every classical branch at prefix $\ell$ has quantum dimension at most
$m\le D$. Put $a=m/D$ and $k=\lfloor1/a\rfloor$. Then
\begin{equation}
 a\sum_{j=0}^{k-1}\lambda_j+(1-ka)\lambda_k\le e_\epsilon.
 \label{eq:profile}
\end{equation}
In particular, if $L=D/m$ is an integer,
\begin{equation}
 \frac1L\sum_{j=0}^{L-1}\lambda_j\le3\epsilon(1-\epsilon).
 \label{eq:average}
\end{equation}
Let $n_\ell(u)=\rank\mathbf1_{[0,u)}(H_\ell)$ and
$Z_\ell(t)=\Tr e^{-tH_\ell}$. For every $u,t>0$,
\begin{align}
 m&\ge\frac{D}{n_\ell(u)}[1-e_\epsilon/u]_+,\\
 m&\ge\frac{D e^{-t e_\epsilon}}{Z_\ell(t)}.
 \label{eq:heat}
\end{align}
\end{theorem}

\begin{proof}
Every composite Kraus matrix through a branch of dimension at most $m$ has
rank at most $m$. Thus $Z$ has Schmidt number at most $m$. For a rank-$m$
vectorized matrix, Schmidt-basis Cauchy--Schwarz gives
\[
 |K\rangle\!\rangle\langle\!\langle K|
 \le mI\otimes\Tr_{\rm left}(|K\rangle\!\rangle\langle\!\langle K|).
\]
Summing and using the uniform marginal gives $0\le Z\le(m/D)I$ and
$\Tr Z=1$. The minimum energy under this density cap is obtained by filling
the lowest eigenspaces of $H_\ell$ first, proving \eqref{eq:profile}.
The mode-count bound follows by charging energy at least $u$ outside the
low-energy subspace. Finally Jensen gives
$\Tr(Ze^{-tH_\ell})\ge e^{-tE}$, while the density cap bounds it above by
$(m/D)Z_\ell(t)$.
\end{proof}

The density-cap relaxation drops the remaining marginal and Schmidt-number
constraints. It is a necessary condition, not a sufficient compression test.
With $\kappa_\ell=(8/3)\lambda_1>0$, a simpler consequence is
\[
 \frac mD\ge1-\frac{8\epsilon(1-\epsilon)}{\kappa_\ell}.
\]
A small minimum gap can make this vacuous even when the full spectral profile
could still give a strong bound.

For variable branch capacity, let $p_y$ be the probability under the maximally
mixed bath input. This is the branch law for the physical Bell-input test,
since the source-plus-bath marginal is exactly tracial. The Choi witness splits
as $Z=\sum_y Z_y$, where $\Tr Z_y=p_y$ and $Z_y$ has Schmidt number at most
$m_y$. The sum of low-capacity branches has marginal at most $I/D$, so the
same reduction argument proves
\begin{align}
 \Pr(m_y>m)&\ge[1-e_\epsilon/u-m n_\ell(u)/D]_+,\\
 \Pr(m_y>m)&\ge[e^{-t e_\epsilon}-(m/D)Z_\ell(t)]_+.
 \label{eq:tail}
\end{align}
If $q_y$ is the number of coherent qubits, sum the bounds at $m=2^j$,
$j=0,\ldots,\ell-1$, to bound $\mathbb E[\min(q_y,\ell)]$. Summing over
fixed physical time intervals gives a qubit-time lower bound. This does not
add the maximum capacities of mutually exclusive histories.

\subsection{What is known about the spectral calculation}

Every moment of $T_\ell T_\ell^*$ has an exact finite transfer representation.
In the normalized Pauli basis let $G$ be the real orthogonal $16\times16$
matrix of $\Ad_{cI+is\mathrm{SWAP}}$, and write $G^{qp}$ for its $4\times4$
source block at bath input $p$ and output $q$. Then
\[
 (T_\ell)_{\boldsymbol q,\boldsymbol p}
 =\langle0|G^{q_\ell p_\ell}\cdots G^{q_1p_1}|0\rangle.
\]
For integer $k\ge1$, summing the cyclic physical indices in
$\mu_k(\ell)=\Tr[(T_\ell T_\ell^*)^k]$ produces a $16^k\times16^k$
transfer matrix $M_k$ and
\[
 \mu_k(\ell)=\langle0|^{\otimes2k}M_k^\ell|0\rangle^{\otimes2k},
 \qquad
 Z_\ell(t)=e^{-t}\sum_{k\ge0}\frac{t^k}{k!}\mu_k(\ell).
\]
The complete index formula and executable checks accompany this report.
The first moment is particularly simple:
\[
 \mu_1(\ell)=4^{\ell-1}(1+3c^{4\ell}).
\]
There is also an exact second moment. With
$u=c^{2\ell}$, $v=r^\ell$, $w=[c^2(c^2-s^2)]^\ell$, and
$h=\ell c^{2\ell-1}s^3$,
\[
 \frac{\mu_2(\ell)}{4^\ell}
 =\frac{12h^2+9u^2+4v^2-6v+5w^2+4}{16}.
\]
It follows either by the finite transfer contraction or by taking the Choi
purity of the complementary two-pass bath channel. The accompanying note
gives the latter derivation and finite checks. These first two moments do
not establish the small-energy mode count needed in the memory bound.
For each fixed $0<\theta<\pi/2$, they do prove convergence of the empirical
spectrum of $H_\ell$ to $\tfrac14\delta_0+\tfrac34\delta_1$ as
$\ell\to\infty$: the normalized mean of $A=T_\ell T_\ell^*$ tends to
$1/4$ and that of $A(I-A)$ tends to zero. This describes a large near-zero
band, not an exact finite-prefix kernel. Its width and extreme lower tail
are still needed for the memory theorem.
This finite representation alone is not an efficient solution of the large
prefix problem. The bath channel has Choi rank at most four, which gives
$\mu_k(\ell)\ge4^\ell/4^k$. Raw moment bounds capable of reducing the
slow-mode count below $2^\ell$ therefore require $k=\Omega(\ell)$.

Analytic trial operators also show why a uniform spectral-gap assertion would
be unsafe. If a normalized bath operator is supported on the first $j$ sites,
then appending identities gives the exact loss identity
\[
 1-\norm{T_\ell(A\otimes I)}_{2,\rm norm}^2
 =r^{\ell-j}\bigl(1-\norm{T_j(A)}_{2,\rm norm}^2\bigr).
\]
Such trials give upper bounds on the lowest eigenvalue on the traceless
subspace, which may be zero, not the lower bounds required by the converse.
Numerical gaps through eight carriers
suggest slow collective modes, but do not establish their large-$\ell$
density. The integrable transfer matrices of Ref.~\cite{claeys} have relevant
eigenvalue and Jordan structure; the required object here is a singular
spectrum. That distinction cannot be suppressed.

\section{Which Schur coherences are already known to matter}

Efficient Schur and Clebsch--Gordan circuits are established
machinery~\cite{bacon}. Their existence does not establish which registers may
be discarded for this source.

For two emitted carriers, put $B_2^a=\tfrac12\Tr_S J_a$. Direct conjugation gives
\begin{equation}
 B_2^a=s^2\left[\sigma_a^{(1)}+c^2\sigma_a^{(2)}
       +cs(\boldsymbol\sigma^{(1)}\times\boldsymbol\sigma^{(2)})_a\right].
 \label{eq:B2}
\end{equation}
In the singlet/triplet-zero subspace, $B_2^z$ has the off-diagonal block
$s^2(s^2X+2csY)$, up to the stated basis convention. Emission order is present
in these operators. The channel is covariant under simultaneous SU(2)
rotations but is not invariant under bath permutations.

\begin{theorem}[Measured-label Schur boundary]
For any SU(2)-covariant qubit-to-$N$-qubit channel whose maximally mixed input
has maximally mixed output, consider the compression that measures total spin
$J$, retains the irrep register $Q_J$, and discards its permutation
multiplicity. Every subsequent decoder has
\[
 \epsilon\ge\frac34\left[1-\sum_{J>0}w_J\frac{J+1}{3J}\right],
\]
where
\[
 w_J=\frac{(2J+1)}{2^N}
 \left[\binom N{N/2-J}-\binom N{N/2-J-1}\right].
\]
As $N\to\infty$, the bound is
$\epsilon\ge\tfrac12-1/\sqrt{2\pi N}+O(N^{-1})$.
\end{theorem}

\begin{proof}
Covariance makes $w_J$ independent of the input and fixes the conditional
qubit-to-spin-$J$ state to
$(I+\alpha_J\boldsymbol r\cdot\boldsymbol J)/(2J+1)$.
Positivity implies $|\alpha_J|\le1/J$. Twirl the final decoder; its vector
part is bounded by the usual extremal spin-$J$ to spin-$1/2$ covariant maps.
Using $\Tr(J_z^2)/(2J+1)=J(J+1)/3$, its recovered visibility is at most
$(J+1)/(3J)$. The $J=0$ branch has no vector signal. Weighting the branches
gives the displayed depolarizing visibility bound and therefore its
half-diamond error. The binomial weights give the stated asymptotic.
\end{proof}

At $N=43\,462$ the finite bound is approximately $0.498074877$.
This theorem explicitly includes measurement of $J$. It does not exclude
coherent retention of adjacent $J$ sectors, tailored multiplicity rotations,
or another causal representation. In fact
$Q_{K-1/2}\oplus Q_{K+1/2}\cong\mathbb C^2\otimes Q_K$ can contain a
logical qubit in the coherence between the two sectors.

\subsection{Making consecutive-pair fusion labels classical}

There is a sharper source-specific obstruction to a local version of that
discard. Measure the singlet/triplet label of each consecutive emitted pair,
but keep both representation blocks, every outcome, and all other available
quantum data. Write $x=s^2$ and
\[
 b=\frac{x^2(2-x)^2}{2},\qquad
 h=(1-x)^2(1+2x)=1-3x^2+2x^3.
\]
After $K$ pairs the exact acquired Pauli Gram coefficient is
\[
 p_K=b\frac{1-h^K}{1-h}.
\]
To prove this, use normalized Hilbert--Schmidt inner products. The new
source-identity component is
$[x(2-x)/2](\sigma_a^{(1)}+\sigma_a^{(2)})$, with squared norm $b$.
The mutually orthogonal source-traceless Pauli-weight groups have squared
norms $(1-x)^4$, $4(1-x)^3x$, and $3(1-x)^2x^2$, whose sum is $h$.
Covariance and fresh-pair orthogonality make the same scalar recurrence valid
with arbitrary old-output operator coefficients.

Weighted Cauchy--Schwarz bounds every final decoder by
\[
 \epsilon_{\rm opt}\ge\tfrac34(1-\sqrt{p_K}),\qquad
 p_K\le\frac{(2-x)^2}{2(3-2x)}\le\frac23.
\]
Consequently the error floor is at least
$\tfrac34(1-\sqrt{2/3})=0.1376275643\ldots$.
The discarded pair-sector coherence matters at fourth order, not at an order
negligible over the acquisition interval. This result does not exclude an
adaptive rotation before a branch measurement.

\section{The folded two-qubit reverse decoder is not a causal update}

A tempting representation initializes an accessible Bell pair $R,A$ and
folds the reverse decoder across the Bell state. The local crossed tensor is
\[
 V_{AE}=(U^\dagger)^{T_S}=cI-2is\Phi_{AE}.
\]
Chronological products of these tensors on the Bell-initialized subspace
represent the exact reverse-dummy channel. The complete prefix map is an
isometry on that subspace. The local update on the retained mixed state,
however, is completely positive but generally not trace preserving.

This failure can be proved operationally, without relying on a failed local
unitary ansatz. For the virtual state after $n$ steps define
\[
 u=c^{2n},\quad v=r^n,\quad
 w=[c^2(c^2-s^2)]^n,\quad h=n c^{2n-1}s^3.
\]
The baseline is $\tau_S\otimes[u\Phi_{RA}+(1-u)I_{RA}/4]$.
The full exact Pauli family is included in the source note. If $D_n$ is the
trace distance between its two states for source inputs $|0\rangle$ and
$|1\rangle$, then
\begin{equation}
 \left|D_n-\frac{1+3|1-4v/3|}{4}\right|
 \le3u+\frac{10}{3}|w|+4|h|.
 \label{eq:Dvirtual}
\end{equation}
In the quartic weak-angle limit $n\theta^4\to a>0$, the right side vanishes,
while the limiting distinguishability first falls to $1/4$ at $e^{-a}=3/4$
and then rises to one.

For a physical causal receiver, the distinguishability of the joint actual
source plus all retained quantum and classical memory cannot increase:
adjoining a fresh tracial carrier and applying the source unitary preserve it,
and the receiver CPTP operation contracts it. Thus this virtual statistic
cannot be a physical causal sufficient statistic. If it is the complete joint
state of the actual source and retained memory at a cut, with no additional
payload-dependent quantum or classical record, even an unrestricted later
continuation obeys
\[
 \epsilon\ge(1-D_n)/2.
\]
At $\theta=0.1$, $n=2\,896$, the right side is at least
$0.37499725066865$. If the two actual states are each within half-trace
distance $\eta$ of the virtual targets, the bound decreases by at most
$\eta$. This is a robust obstruction to this particular statistic; it is
not a two-qubit-memory lower bound for all possible statistics.

Postselection does not repair the deterministic bill. Normalizing $V$ gives
$K=V/\sqrt{1+3s^2}$. The all-success event has probability
$(1+3s^2)^{-N}$, independent of the payload. At $\theta=0.1$, $N=43\,462$,
this is approximately $10^{-556.10}$. The input is one unknown specimen and
cannot be repeatedly prepared after a failed attempt.

\section{Purified-dummy blocks: an exact collective-decoder obstruction}

A more substantial physical proposal stores each block of $b$ outputs,
prepares a fresh accessible Bell pair $R,A$, reverse-decodes the block using
$R$, discards the block, and retains \emph{both} $R$ and $A$. The accessible
purifier is prepared at the receiver; it is not a bath preparation key.
After $K$ blocks, any collective decoder of all $2K$ retained qubits is
allowed. The proposed peak storage would be $b+2N/b$ plus workspace.

Let $u,v,w,h$ have the preceding definitions with $n=b$, and put
$a=1-v$, $d=\sqrt{(1-u)(1+3u)}$. The baseline is
$\sigma=\tau_S\otimes\omega_{RA}$, with
$\omega_{RA}=u\Phi_{RA}+(1-u)I/4$. In the symmetric Petz metric
\[
 \langle X,Y\rangle_\sigma
 =\Tr(X\sigma^{-1/2}Y\sigma^{-1/2}),
\]
the source-identity and source-traceless parts of the Pauli signal are
orthogonal. Their exact squared norms are
\begin{align}
 p_b={}&\frac{a^2}{2(1-u)}+\frac{a^2+4h^2}{2d},\\
 q_b={}&\frac{(6u+v+5w)^2}{36(1+3u)}\nonumber\\
 &+\frac{36h^2+27u^2-24uv-30uw+7v^2+10vw+10w^2}{36(1-u)}\nonumber\\
 &+\frac{36h^2+9u^2+12uv-30uw+4v^2-20vw+25w^2}{36d}.
 \label{eq:blockpq}
\end{align}
These are obtained by splitting $\sigma$ into the Bell and orthogonal
eigenspaces and evaluating their three quadratic blocks. The RLD quantity
$\Tr(X^2\sigma^{-1})$ cannot be substituted for this metric.

The exact repeated-block law is
\begin{equation}
 p_K=p_b\frac{1-q_b^K}{1-q_b},\qquad q_K=q_b^K.
 \label{eq:blockrec}
\end{equation}
Baseline block outputs are independent, although the actual payload-bearing
outputs need not factor. The metric weights factor, covariance makes the
three source-Pauli Gram matrices scalar, and fresh trace-zero export is
orthogonal to old exports. This proves \eqref{eq:blockrec} even with all
old-output correlations retained.

The collective Petz decoder has visibility $p_K$. Every collective decoder
satisfies
\begin{equation}
 \epsilon\ge\tfrac34(1-\sqrt{p_K}).
 \label{eq:blocklower}
\end{equation}
Indeed, for each recovered Pauli the adjoint observable is a Hermitian
contraction. Symmetric weighted Cauchy--Schwarz bounds its overlap with the
input Pauli by $\sqrt{p_K}$; summing the three overlaps bounds entanglement
fidelity by $(1+3\sqrt{p_K})/4$.

For every $b(\theta)=o(\theta^{-4})$,
\begin{equation}
 \frac{p_b}{1-q_b}\longrightarrow0\qquad(\theta\to0).
 \label{eq:blocklimit}
\end{equation}
A complete proof in the accompanying note separates $b\theta^2\to0$,
$b\theta^2\to x\in(0,\infty)$, and $b\theta^2\to\infty$.
In the first regime $p_b\sim b^{3/2}\theta^5$ and
$1-q_b\sim2b\theta^2$. In the second $p_b=O(\theta^2)$ and $q_b$ tends
to a value strictly below one. In the third $p_b\to0$, $q_b\to1/3$.
At $b\theta^4\to\alpha>0$,
\[
 \frac{p_b}{1-q_b}\longrightarrow
 \frac{(1-e^{-\alpha})^2}{1-e^{-2\alpha}/3}.
\]
High accuracy still requires a block on the quartic logarithmic scale.
For unequal predetermined blocks, the total acquired coefficient is bounded
by $\max_b p_b/(1-q_b)$ over the used blocks.

\begin{center}
\begin{tabular}{r r}
\toprule
Block length $b$, at $\theta=0.1$ & $p_b/(1-q_b)$\\
\midrule
1 & 0.000583746\\
100 & 0.004314638\\
1\,000 & 0.012305302\\
3\,000 & 0.081359214\\
10\,000 & 0.415496571\\
40\,000 & 0.962856321\\
\bottomrule
\end{tabular}
\end{center}
Even infinitely many $40\,000$-carrier blocks in this architecture have
error at least $0.01406068$. The proposed square-root-memory block interface
therefore fails on information retention, independently of gate compilation.
Other block instruments and other retained block modes remain outside this
theorem.

\section{A broader one-qubit boundary and constructive candidates}

There is a further proved boundary beyond scalar partial-SWAP receiver
updates. Retain one tracial qubit $M$ and allow an arbitrary time-dependent,
nonselective SU(2)-covariant CPTP map $ME\to M$ at each arrival. This includes
dissipative channels implemented using initialized ancillas and reset. It
does not include an extra retained measurement history that steers the source
conditionally, or an extra quantum register.

Such a channel has Heisenberg vector action
\[
 \Phi^*(\boldsymbol\sigma)=\alpha\boldsymbol\sigma_M
 +\beta\boldsymbol\sigma_E
 +\gamma(\boldsymbol\sigma_M\times\boldsymbol\sigma_E).
\]
Its exact complete-positivity region is
\[
 -\tfrac12\le z=\alpha+\beta\le1,\qquad
 (\alpha-\beta)^2+4\gamma^2\le(1+2z)/3.
\]
The source-plus-memory signal closes in the three vector components
$\boldsymbol\sigma_S$, $\boldsymbol\sigma_M$, and
$\boldsymbol\sigma_S\times\boldsymbol\sigma_M$.
Using this CP region gives, uniformly over the number of collisions and all
such update sequences, the recovered memory visibility bound
\[
 |B_N|\le\sin^2\theta\left(1+\frac{8}{3-\cos\theta}\right).
\]
Consequently every final decoder obeys
\[
 \epsilon\ge\frac34\left[1-\sin^2\theta
                 \left(1+\frac8{3-\cos\theta}\right)\right]_+.
\]
At $\theta=0.1$ this lower bound is $0.71269934$. The optimal visibility
in this restricted class is $\Theta(\theta^2)$: the upper bound has that
scale, and retaining just the first emitted carrier attains visibility
$\sin^2\theta$. This lower example does not perform a full SWAP at the
source. The complete CP-region and Lyapunov proof is in the accompanying
one-qubit note; 295 direct finite channel checks verify its recurrence.

\subsection{Growing-spin receivers}

The accompanying files specify a physical growing-spin receiver, not merely
a formal tensor contraction. A qubit $R$ and spin-$J$ register are retained;
each incoming carrier undergoes a covariant control and a permitted
Clebsch--Gordan branch instrument. The encoded quantum dimension is
$O(N)$, so the register uses $O(\log N)$ qubits, with binary arithmetic
workspace. This is a memory bill for the instrument, not a recovery theorem.

For the tracial conditional family, the invariant joint Choi state of
reference, source, receiver qubit and spin-$J$ has only 20 Hermitian-block
parameters for $J\ge3/2$. The next carrier temporarily raises this to 70.
General covariant controls reduce to constant-size multiplicity blocks.
The exact immediate-loss-minimizing branch choice reduces to a three-variable
quadratic optimization on the unit sphere. Its executable implementation
keeps the full realized classical history and the permitted coherences.

That particular adaptive policy performs poorly. At $\theta=0.1$ after
$5\,000$ collisions, three recorded histories have acquired Gram
coefficients approximately $0.000878$, $0.000961$, and $0.004202$.
Their remaining source-traceless signal is already very small. This is a
diagnostic of a specified local control objective. It proves no impossibility
for other adaptive policies, nontracial initialized workspace, or a coherent
multi-irrep receiver.

The source notes also examine a full finite rotor, retain its left/right
multiplicity and inter-irrep coherence, and derive an optimal final readout
for the relevant finite spin family. No high-fidelity rotor reader is claimed.
The pure-singlet, nontracial variant was implemented with its actual weighted
Choi state. In the tested preservation policy the pure resource becomes
lossless storage of the first two outputs; a recorded trajectory at
$\theta=0.3$, $N=600$ ends with optimal conditional final-readout entanglement
fidelity $0.48206991236$. This is one implemented policy and history, not a
bound on all pure initialized receivers or a statistically certified policy
average. Coherent sparse-flag variants remain unproved; a small excitation
count alone is not a recovery theorem.

\subsection{Established alternatives and why they do not finish the task}

Bacon--Chuang--Harrow provide efficient Schur and Clebsch--Gordan circuitry,
with accuracy-dependent compilation cost~\cite{bacon}. The new requirement
here is to prove that the discarded register is irrelevant to the arbitrary
payload. Equations \eqref{eq:B2} and the pair-loss theorem show concrete
failures of that premise.

Blume-Kohout--Croke--Zwolak give sequential coherent gathering for a finite
set of pure states and matrix product states of bounded Schmidt
rank~\cite{bcz}. Their mixed-state extension exploits permutation symmetry
of identically prepared copies. Our bath average is maximally mixed, the
payload-bearing output is ordered and correlated, and its physical
purifiers are unavailable. A small matrix-product operator bond or a compact
formula for three Pauli observables is not the pure-subspace condition of
that construction.

Homogenization and recovery-existence work already relates partial-SWAP
reservoirs to correctability~\cite{yosifov}. The accepted warm package supplies
the stronger explicit all-output operation \eqref{eq:baseline}. Neither a
correctability criterion nor source forgetting furnishes the required
forward causal compression.

The integrable transfer-matrix construction~\cite{claeys} supplies useful
spectral and representation machinery. Its homogeneous Jordan blocks mean
that eigenvalues cannot simply replace the singular spectrum in
\eqref{eq:H}. A pictured double row is also not a proof of an applicable
reflection-equation solution; the exact boundary must be checked.

\section{Resources and finite errors}

For the accepted all-output reader, the joint scaling is explicit:
\[
 Q_{\rm retained}=N_{\rm all}+O(1),\qquad
 G_{\rm inverse}=N_{\rm all},\qquad
 T_{\rm collect}=N_{\rm all}\Delta.
\]
Here the collection formula assumes regular arrivals with actual encounter
spacing $\Delta$. A source Hamiltonian
with norm budget $\Omega_S$ needs encounter action at least $|\theta|$, hence
duration at least $|\theta|/\Omega_S$ for the stated pulse convention.
Routing cost depends on the delivered physical geometry. A fixed-port serial
register can require quadratic SWAP traffic; direct addressing or moving
carriers is an additional layout premise whose transport, timing, and
protection must be paid. Linear abstract inverse-gate count does not erase
that premise.

For the concrete nearest-neighbor junction layout in the accepted physical
package, capture uses $N(N-1)/2$ routing SWAPs and decoding uses $N-1$.
Its total source, inverse, and routing gate count is
\[
 G_{\rm junction}=\frac{N^2+5N-2}{2}.
\]
If $\tau_s,\tau_d,\tau_w$ are source, inverse, and full-SWAP durations,
and $\tau_{\rm ctl}$ bounds control latency per gate, then
\begin{align*}
 T_c&=N\tau_s+\frac{N(N-1)}2\tau_w
             +\left(N+\frac{N(N-1)}2\right)\tau_{\rm ctl},\\
 T_d&=N\tau_d+(N-1)\tau_w+(2N-1)\tau_{\rm ctl}.
\end{align*}
The supplied conservative memory exposure is
$(N+1)(T_c+T_d)$, with additional preparation and waiting exposure if present.
Thus this particular bounded-control layout has quadratic operation time and
cubic memory space-time at fixed gate durations. At the numerical point it
uses $944\,581\,376$ two-qubit gates. No routing optimum is asserted.

Even before decoding, storing all outputs incurs
$N(N-1)\Delta/2$ qubit-time. At the numerical point this is
$944\,450\,991\Delta$. Memory protection is therefore a substantive cost.
No factory or large protected receiver was simulated in this work.

Let $\eta_{
m prep}$, $\eta_{
m src}$, $\eta_{
m rec}$,
$\eta_{
m route}$ and $\eta_{
m store}$ be sums of half-diamond
implementation errors in the corresponding operations, with state
preparation errors measured in half trace norm. The ordinary hybrid bound
gives
\[
 \epsilon_{
m actual}\le\epsilon_{
m ideal}
 +\eta_{
m prep}+\eta_{
m src}+\eta_{
m rec}
 +\eta_{
m route}+\eta_{
m store}.
\]
Systematic angle errors do not average away. A conservative allocation is
absolute source and inverse-angle error $O(\epsilon/N)$ per encounter,
with independent budgets for routing, memory noise, and initialization.
Calibrating to that precision consumes actual known preparations and
measurements; the theorem does not certify unobserved independence or absent
couplings from a short calibration record.

The accepted finite factory and qualification construction is preserved in
full. An ideal independent Pauli-randomization factory uses two random bits
and at most two single-qubit Pauli pulses per tracial register, with every
operating key copy inaccessible. Its general sufficient independent-factory
qualification cost is $O(N^2\epsilon_p^{-2}\log(N/\alpha_p))$ trials.
For the junction circuit, separate fixed-dimensional gate qualification at
error $O(\epsilon_g/G_{\rm junction})$ has the conservative cost
$O(N^5\epsilon_g^{-2}\log(N/\alpha_g))$ trials. These bills assume the
stated stationary finite-dimensional model, trusted preparations/readout,
and actual independence where used. They are not optimized or demonstrated
hardware costs. Protected implementations must add their encoding, factory,
reset, routing, latency, and certified logical-error costs.

For binary-encoded spin receivers, a compiled update of depth $G$ at local
clock rate $\Omega_R$ needs service time $\gtrsim G/\Omega_R$. An input
spacing faster than that creates a quantum backlog. It must be included in
memory; one cannot claim logarithmic storage while assuming infinitely fast
arithmetic. The mathematically specified spin instruments have
$O(N\,\mathrm{polylog}(N/\delta))$ elementary processing at precision
$\delta$ under the corresponding service-time premise. Their missing
high-fidelity recovery guarantee remains decisive.

A direct full SWAP transfers the arbitrary input and reference in one
encounter if the source interface permits action $\pi/2$. It remains the
simpler baseline in that interface. The fixed weak encounter here does not
supply it, and receiver workspace does not lengthen an already completed
source encounter.

\section{Evidence, preservation, and the exact remaining gap}

The computational checks are finite matrix identities and policy diagnostics,
not physical experiments. The crossed-state closed form was checked in 100
finite cases with maximum residual $1.33\times10^{-15}$. The purified-block
weighted matrices agree with the formulas to $6.31\times10^{-14}$; direct
two-block channels agree with the scalar recurrence to
$2.36\times10^{-14}$. Pair-fusion identities were checked against full
matrices at $N=2,4,6$ with residual at most $1.2\times10^{-15}$.
The archive records numerical parameters, scripts, and output receipts.
None of these checks is presented as simulation of a $43\,462$-carrier
quantum register.

The complete earlier finite-error payload-memory theorem, its short-record
construction, prefix and occupancy statements, noisy-emission correction,
and all three companions are preserved without modification in the included
warm archive. The tracial acquisition package is also retained unchanged.
The programmed emitter and early-record decoder remain separate from the
present homogeneous weak source. Their resource and timing premises are not
silently transferred.

There are two precise unresolved endpoints:
\begin{enumerate}
\item A positive solution must give a CPTP causal instrument on the actual
once-through carriers, prove reference-uniform recovery at
$N=\Theta(\theta^{-4}\log(1/\epsilon))$, and bound all quantum registers
and their service backlog jointly in $\theta,\epsilon,N$.
\item A negative solution must evaluate a genuinely general converse, such
as \eqref{eq:profile}--\eqref{eq:tail}, strongly enough on this family to
prove the required memory boundary, with a matching operation. A failing
spin update or block decoder does not establish it.
\end{enumerate}
This return supplies neither endpoint. It supplies a source-specific general
converse and closes several tempting but invalid reductions, with their
proofs and execution costs exposed. The optimal memory/processing tradeoff
and the requested compressed arbitrary-input acquisition operation remain
unresolved.

\begingroup
\footnotesize
\begin{thebibliography}{9}
\setlength{\itemsep}{1pt}
\bibitem{bacon} D. Bacon, I. L. Chuang, A. W. Harrow,
\emph{Efficient Quantum Circuits for Schur and Clebsch--Gordan Transforms},
Phys. Rev. Lett. 97, 170502 (2006).
\url{https://arxiv.org/abs/quant-ph/0407082}.
\bibitem{bcz} R. Blume-Kohout, S. Croke, M. Zwolak,
\emph{Ideal state discrimination with an $O(1)$-qubit quantum computer}
(2012). \url{https://arxiv.org/abs/1201.6625}.
\bibitem{claeys} P. W. Claeys, J. Herzog-Arbeitman, A. Lamacraft,
\emph{Correlations and commuting transfer matrices in integrable unitary
circuits}, SciPost Phys. 12, 007 (2022).
\url{https://arxiv.org/abs/2106.00640}.
\bibitem{yosifov} A. Yosifov, A. Iyer, D. Ebler, V. Vedral,
\emph{Quantum Homogenization as a Quantum Steady State Protocol on NISQ
Hardware}, Phys. Rev. A 109, 032624 (2024).
\url{https://arxiv.org/abs/2412.14544}.
\end{thebibliography}
\endgroup
\end{document}
